\chapter{多源脑电信号生成式建模理论基础}
\label{chapter2}

本章从理论角度系统构建多源脑电信号生成式建模的数学基础。首先给出多源脑电信号的形式化表示，明确其异构性与多层级信息结构，为后续建模提供统一的符号体系与概率论框架。随后分析传统判别式范式在表达能力、数据效率与泛化能力等方面的理论局限，并据此引出生成式建模的核心优势，包括连续表达能力、自监督学习潜力、不确定性建模与通用表征学习等。最后提出信号级--表征级的双层级生成式建模框架，分别给出可操作的目标函数与关键设计要点，并明确各层之间的功能分工与信息衔接关系。该理论分析揭示了生成式建模相比判别式方法在多源EEG信号处理中的优势，为后续方法与实验奠定统一且可扩展的理论框架。


\section{多源脑电信号的形式化表示}
\label{sec:formal}
多源EEG信号作为本研究的核心数据对象，具有来源多样、结构复杂、统计机制异质与信息层级丰富等特点。从信息论的角度来看，EEG承载大脑认知活动的高维时--空--频信息，并通过多路体积传导与复杂的非线性机制编码于观测信号之中。因此，有必要同时建立面向观测的张量化表示与面向生成的概率化表示，并用信息结构统一组织不同尺度的统计属性。三者互为补充：张量视角回答“数据长什么样”；概率视角回答“数据作为随机变量如何生成”；信息结构回答“哪些统计与内在依赖应被保留”。

\subsection{多源脑电信号的基本定义}
为清晰刻画多源EEG的外在形态与内在随机性，本节先从确定性张量结构给出数据形状与元信息的描述，再在此基础上引入概率测度空间以体现随机本质。多源EEG信号来源于不同被试、不同采集设备、不同实验范式或不同时间段的观测集合，见定义~\ref{def:eeg-tensor}。

\begin{definition}[多源EEG信号的张量化表示]
\label{def:eeg-tensor}
多源EEG可表示为异构时--频--空张量集合
\begin{equation}
\label{eq:tensor-collection}
\mathcal{X} = \Big\{ X^{(s)} \in \mathbb{R}^{C_s \times T_s \times F_s} \Big\}_{s=1}^{S},
\end{equation}
式中，$S \in \mathbb{N}^+$为来源数量，$X^{(s)}$为第$s$个来源的EEG张量，$C_s, T_s, F_s \in \mathbb{N}^+$分别为该来源的通道数、时间采样点数与频率分量数。每个来源配套元数据$M^{(s)} \in \mathbb{R}^{d_m}$（其中$d_m \in \mathbb{N}^+$），记录电极布局、参考电极、设备增益与频响、采样率、阻抗、被试生理信息与任务条件等上下文。为兼容不规则采样与缺失，增加二值掩码$R^{(s)} \in \{0,1\}^{C_s \times T_s \times F_s}$表示可用条目，其中$R^{(s)}_{c,t,f} = 1$表示位置$(c,t,f)$的观测值有效。
\end{definition}

张量化表示明确了数据的几何形状与跨源异构特性，为理解数据的基本结构提供了重要基础。时间维度$T_s$反映了信号的动态演化过程，不同来源的采样率差异导致时间分辨率的不一致。通道维度$C_s$体现了空间采样的复杂性，不同设备的电极配置差异（如64导联vs 128导联系统）导致空间分辨率和覆盖范围的显著变化。频率维度$F_s$则刻画了信号的频谱特性，不同的滤波设置和频谱分析方法会影响频域表示的精度和范围。元数据$M^{(s)}$的引入使得模型能够显式地利用采集过程中的先验信息，为后续的标准化和对齐提供依据。掩码$R^{(s)}$的设计则考虑了实际应用中常见的数据质量问题，如电极脱落、信号饱和、伪影污染等导致的数据缺失。张量化表示仅描述了数据的几何形状，未触及生成机制和统计性质。为此引入概率测度空间，以支持统计推断、似然建模与不确定性量化，见定义~\ref{def:eeg-prob}。


\begin{definition}[多源EEG的概率空间表述]
\label{def:eeg-prob}
对于第$s$个来源，设$(\Omega_s, \mathcal{F}_s, P_s)$为完备概率空间，其中$\Omega_s$为样本空间，$\mathcal{F}_s$为$\sigma$-代数，$P_s$为概率测度。定义随机张量$\mathbf{X}^{(s)}: \Omega_s \to \mathbb{R}^{C_s \times T_s \times F_s}$为$\mathcal{F}_s$-可测函数。多源EEG的概率描述为
\begin{equation}
\label{eq:prob-collection}
\mathcal{P} = \Big\{ \big(\mathbf{X}^{(s)}, (\Omega_s, \mathcal{F}_s, P_s)\big) \Big\}_{s=1}^{S},
\end{equation}
式中，$\mathbf{X}^{(s)}$为第$s$个来源的随机张量，$(\Omega_s, \mathcal{F}_s, P_s)$为对应的概率空间，$P_s$反映设备噪声、个体差异、任务操控与环境扰动等因素导致的统计变异。进一步地，可在统一潜在源空间引入线性观测算子$A_s \in \mathbb{R}^{C_s \times N_s}$与加性噪声$\boldsymbol{\epsilon}_s(t,f) \sim \mathcal{N}(0, \Sigma_s)$，建立观测模型
\begin{equation}
\label{eq:observation-model}
\mathbf{X}^{(s)}(c,t,f) = \sum_{n=1}^{N_s} A_s(c,n) \mathbf{S}_n(t,f) + \boldsymbol{\epsilon}_s(c,t,f),
\end{equation}
式中，$\mathbf{S}_n(t,f) \in \mathbb{R}$表示第$n$个源在时刻$t$、频率$f$处的活动强度，$N_s \in \mathbb{N}^+$为源的数量，$A_s(c,n)$为体积传导算子的第$(c,n)$个元素，$\boldsymbol{\epsilon}_s(c,t,f)$为第$s$个来源在通道$c$、时刻$t$、频率$f$处的加性噪声，$\Sigma_s$为噪声协方差矩阵。
\end{definition}

定义~\ref{def:eeg-prob}中的观测模型提供了理解多源差异的统一框架。共享的皮层源活动$\mathbf{S}_n(t,f)$体现了大脑神经活动在时频域的本质特征，这些特征在不同个体和条件下具有一定的共性，构成了跨源信息传递的生理基础。观测算子$A_s$刻画了从源空间到传感器空间的线性映射关系，该映射主要由体积传导效应决定，受到头部几何形状、组织电导率分布、电极空间配置等解剖学和技术因素的影响。不同来源的$A_s$差异正是导致多源异构性的核心原因，体现了个体差异、设备差异和实验条件差异在观测层面的综合表现。加性噪声$\boldsymbol{\epsilon}_s(c,t,f)$综合反映了各种随机扰动因素的影响，包括设备热噪声、量化噪声、生理伪影（如眼电、肌电干扰）、环境电磁干扰等，其协方差结构$\Sigma_s$编码了不同通道间噪声的相关性特征。

在上述概率框架之上，还需以信息结构组织EEG的多维度、跨尺度统计属性，见定义~\ref{def:info-struct}。
% 概率空间提供统计刻画，信息结构提供机制与统计层面的分解，共同构成建模基础。

\begin{definition}[EEG信号的信息结构]
\label{def:info-struct}
设$X(c,t,f)$为通道$c$、时刻$t$、频率$f$处的信号值。定义信息结构为四个信息分量的集合
\begin{equation}
\label{eq:info-struct}
\mathcal{I}(X) = \big\{ I_{\text{tmp}}(X), I_{\text{spa}}(X), I_{\text{spc}}(X), I_{\text{crs}}(X) \big\},
\end{equation}
式中，$I_{\text{tmp}}(X)$为时间动态信息，描述自回归性、事件相关电位与相位动力学；$I_{\text{spa}}(X)$为空间分布信息，刻画脑区协同、体积传导影响与功能网络拓扑；$I_{\text{spc}}(X)$为频域特性信息，反映功率谱密度、相位同步与跨频耦合；$I_{\text{crs}}(X)$为跨维交互信息，编码时--频、空--时与空--频的耦合项，如时频自回归耦合、传递熵与相干性的时变性。
\end{definition}

时间信息$I_{\text{tmp}}$可由动力系统与时间序列模型刻画事件相关电位与相位动力学，反映神经活动的时序规律和因果关系。在BCI应用中，时间动态信息对于理解用户意图的演化过程和预测未来状态具有重要意义。空间信息$I_{\text{spa}}$可通过空间滤波、源定位与网络连通性分析揭示功能分工，体现了大脑不同区域的协同工作模式。现代神经科学研究表明，认知功能往往需要多个脑区的协调配合，空间信息的准确建模对于理解复杂认知过程至关重要。频谱信息$I_{\text{spc}}$通过功率谱与相位耦合衡量不同频段的神经节律，反映了大脑的振荡活动模式，不同频段的神经振荡与不同的认知功能相关联。交互信息$I_{\text{crs}}$捕捉跨维度的非线性耦合，是理解复杂认知的重要纽带。这种高阶交互信息往往包含了大脑信息处理的关键特征，但也是最难建模和分析的部分。


\subsection{多源脑电信号的异构性分析}
多源EEG的异构性不仅体现在外在形态（通道数、采样率、频带设置等），更深层地体现在统计分布、结构模式与功能编码差异。
这种异构性是多源EEG建模面临的核心挑战，需要在理论框架中给予充分考虑，见定义~\ref{def:hetero}。

\begin{definition}[多源信号的异构性特征]
\label{def:hetero}
对于多源EEG系统$\{(\mathbf{X}^{(s)}, (\Omega_s, \mathcal{F}_s, P_s))\}_{s=1}^{S}$，记第$s$个来源随机张量$\mathbf{X}^{(s)}$的概率分布为$\mathcal{D}_s$，结构算子为$\Phi_s: \mathbb{R}^{C_s \times T_s \times F_s} \to \mathcal{H}$（其中$\mathcal{H}$为结构特征空间），语义映射为$f_s: \mathbb{R}^{C_s \times T_s \times F_s} \to \mathcal{C}$（其中$\mathcal{C}$为认知状态或任务标签空间）。多源EEG的异构性包含以下三个层面：

\begin{equation}
\label{eq:hetero}
\left\{
\begin{aligned}
\text{分布异构：}&\quad \mathcal{D}_{s_1} \neq \mathcal{D}_{s_2}, \quad \forall\, s_1 \neq s_2 \in \{1,2,\ldots,S\};\\
\text{结构异构：}&\quad \Phi_{s_1}\!\big(\mathbf{X}^{(s_1)}\big) \neq \Phi_{s_2}\!\big(\mathbf{X}^{(s_2)}\big), \quad \forall\, s_1 \neq s_2 \in \{1,2,\ldots,S\};\\
\text{语义异构：}&\quad f_{s_1}^{-1}(c) \neq f_{s_2}^{-1}(c), \quad \forall\, c \in \mathcal{C}, \, s_1 \neq s_2 \in \{1,2,\ldots,S\}.
\end{aligned}
\right.
\end{equation}
\end{definition}


分布异构性反映了不同来源信号在统计特性上的根本差异。这种差异源于多个层面的因素：设备层面的差异包括不同厂商设备的频率响应特性、背景噪声水平、动态范围、采样精度、电极材质等技术参数的不同；环境层面的差异包括实验室和自然环境、静音和噪声环境、电磁屏蔽和开放空间等条件的变化；生理层面的差异包括个体间的头皮厚度、脑容量、皮层折叠模式、神经元密度等解剖结构变异，以及被试的动态生理状态变化，如情绪波动（焦虑、兴奋、抑郁）、疲劳状态、药物影响、昼夜节律变化、激素水平波动等生理因素。这些因素共同作用，导致不同来源的信号在均值、方差、偏度、峰度等统计矩上呈现显著差异。数学上，分布异构性可通过各种概率距离度量来量化，如Kullback-Leibler散度衡量分布间的信息差异，Wasserstein距离刻画分布间的几何差异，最大均值差异则提供了非参数的分布比较方法。

结构异构性体现在信号的时空模式、频谱特征和动力学行为的深层差异。大脑作为复杂的动力学系统，其活动模式受到多种因素的调节。个体间的大脑解剖结构差异，如皮层厚度、沟回模式、白质连接等，会导致相同认知任务在不同个体中激活不同的脑区组合和连接模式。任务相关的结构差异则反映了认知策略的多样性，同一任务可能通过多种不同的神经机制来实现。设备相关的结构差异主要来源于不同的空间采样方案和信号处理流程，这些差异会影响信号的空间分辨率和频谱特性。结构算子$\Phi_s(\cdot)$可以包括多种形式的分析方法：连通性分析用于揭示脑区间的功能联系，拓扑特征提取用于刻画网络的几何性质，动力学系统分析用于理解信号的演化规律。这些结构特征的差异使得基于固定结构假设的传统方法难以适应多源数据的复杂性。

语义异构性则表明相同的认知状态在不同个体或条件下可能表现为完全不同的神经电活动模式。这种异构性涉及大脑功能组织的个体差异、认知策略的多样性以及神经可塑性的影响。在神经生理层面，大脑的功能可塑性使得相同的认知功能可能通过不同的神经回路来实现，这种可塑性在发育过程、学习过程和损伤恢复过程中都有体现。不同个体的大脑功能网络拓扑结构存在显著差异，导致相同的认知任务可能激活不同的脑区或采用不同的信息处理路径。在行为认知层面，认知策略的个体差异进一步加剧了语义异构性，例如在运动想象任务中，有些被试倾向于视觉化策略（想象动作的视觉场景），而另一些被试则偏好动觉策略（感受肌肉收缩和关节运动），这些策略差异会在EEG信号中表现为完全不同的时频空间激活模式。在技术实现层面，即使采用相同的实验范式，不同的指导语、任务难度设置、反馈方式等实验细节都可能导致被试产生不同的认知状态，从而形成语义层面的异构性。数学上，语义异构性体现为映射$f_s^{-1}(c)$在不同来源间的差异，即相同的标签$c$对应的EEG模式集合在不同来源中可能完全不重叠，这使得跨个体和跨条件的认知状态解码成为一个极具挑战性的问题，需要开发能够捕捉和适应这种语义差异的建模方法。


三类异构相互交织，功能异构往往通过结构差异体现在观测信号上，结构差异进一步在统计分布层面被放大和复杂化。这种多层次的异构性要求生成式建模方法必须具备足够的灵活性和适应性，能够在统一框架下处理不同层次的差异，同时保持对共性特征的有效学习和利用。


\section{判别式建模理论分析}
判别式建模聚焦条件分布$P(Y\mid X)$的学习，强调任务导向的输入--输出映射。尽管其在特定场景中高效实用，但面对定义~\ref{def:hetero}所刻画的多源异构与复杂统计结构，判别式范式存在内在瓶颈。本节在形式化定义基础上，从表达能力、数据效率与泛化稳健性三个方面进行理论分析，并辅以与EEG场景相关的深入讨论。


\subsection{判别式建模的形式化表述}
判别式建模的核心是学习从观测数据到标签的映射关系。给定单次观测和分类器，映射关系为
\begin{equation}
\label{eq:disc-map}
\hat{y} = f_{\theta}\big(\phi(X)\big),
\end{equation}
式中，$X \in \mathbb{R}^{C \times T}$为单次观测，$\phi:\mathbb{R}^{C \times T}\to \mathbb{R}^{d}$为特征映射，$f_{\theta}:\mathbb{R}^{d}\to \mathcal{Y}$为分类器，$\mathcal{Y}=\{1,\dots,K\}$为离散标签集合，$\theta$为模型参数，$d$为特征维度。给定标注数据集，经验风险最小化为
\begin{equation}
\label{eq:emp-risk}
\theta^{\star} = \arg\min_{\theta}\; \frac{1}{N} \sum_{i=1}^{N} L\!\left(y_i, f_{\theta}\!\big(\phi(X_i)\big)\right),
\end{equation}
式中，$L(\cdot,\cdot)$为损失函数，$N$为样本数。

该过程隐含几个重要假设：首先是分布一致性假设，即训练和测试数据来自相同的分布；其次是样本充分性假设，即有限样本能够代表真实分布的特性；第三是决策边界稳定性假设，即存在稳定的决策边界将不同类别分离。判别式方法的核心思想是在特征空间中学习最优的决策边界，使得不同类别的样本能够被正确分类。
在多源EEG分析中，判别式建模面临的具体挑战包括：特征空间的高维性和稀疏性使得决策边界的学习变得困难；个体间的显著差异导致决策边界的不稳定性；时间非平稳性使得固定的决策边界无法适应信号特性的动态变化；多任务场景下的标签空间不一致性进一步加剧了建模困难。


\subsection{判别式建模的局限性分析}

（一）表达能力的有界性

判别式模型的输出空间为有限集合，天然形成信息瓶颈。即便输入空间连续且复杂，输出也被限制在有限个离散标签之内，造成了严重的信息压缩和表达能力限制。

\begin{theorem}[判别式输出空间的表达瓶颈]
\label{thm:disc-cap}
若输出空间$\mathcal{Y}=\{1,\dots,K\}$有限，则任意判别式映射$f:\mathcal{X}\to \mathcal{Y}$满足$|f(\mathcal{X})|\le K$，从而上界输出信息量为$\log_2 K$比特。
\end{theorem}
\begin{proof}
显然$f(\mathcal{X}) \subseteq \mathcal{Y}$且$|\mathcal{Y}|=K$，故$|f(\mathcal{X})|\le K$，最大香农信息量为$\log_2 K$。
\end{proof}

定理~\ref{thm:disc-cap}指出了判别式建模在表达能力上的根本限制。在EEG-BCI应用中，这种限制表现得尤为突出。人类大脑的神经活动是连续的、多维的、动态变化的，包含了丰富的认知信息和情感状态信息。大脑的信息编码具有高度的连续性和多维性特征。单个神经元的发放模式、神经元群体的协同活动、不同脑区间的动态连接等，都包含着丰富的连续变化信息。然而，将复杂的神经信息压缩到几个离散标签中必然导致大量信息的丢失。这种信息瓶颈不仅限制了系统的功能多样性，也浪费了珍贵的神经信号资源。从信息理论的角度看，这种压缩是不可逆的，一旦信息被丢失，就无法通过后续处理来恢复。


（二）数据利用效率的低下

EEG信号的标注需要专业的神经科学知识、严格的实验控制和大量的时间投入，使得高质量标注数据的获取成本极高。这种数据稀缺性与判别式学习对大量标注数据的依赖形成了突出矛盾。在实际的EEG实验中，一个完整的数据采集过程通常需要数小时的准备时间，包括电极放置、皮肤处理、阻抗调节、实验说明等步骤。数据采集阶段则需要被试在严格控制的环境中完成重复性任务，这个过程不仅耗时（通常需要2--4小时），而且对被试的注意力和配合度要求很高。实验的标准化要求使得数据采集的规模受到很大限制，实验室通常只能获得有限规模的数据集。此外，个体间的神经活动模式存在显著差异，这种个体特异性进一步加剧了数据效率问题。针对特定个体标注的数据往往无法直接迁移到其他个体，严重制约了BCI系统的可扩展性，使得系统部署的成本和时间大大增加。此外，标注质量的不一致性也是影响数据效率的重要因素。EEG信号的标注往往依赖于实验任务的设计和被试的主观报告，这种标注方式容易引入噪声和不一致性。

为量化判别式学习对标注规模的依赖，可以从统计学习理论的角度分析其样本复杂度。

\begin{theorem}[VC维度视角的样本复杂度\cite{vapnik1999overview}]
\label{thm:sample}
设假设类的VC维为$d_{\mathrm{vc}}$，0-1损失，给定误差容忍$\epsilon \in (0,1)$与置信度$1-\delta$，若标注样本数满足式~\eqref{eq:vc-sample}，则以至少$1-\delta$的概率，经验风险最小化得到的假设的泛化误差不超过$\epsilon$：
\begin{equation}
\label{eq:vc-sample}
N \ge \frac{4}{\epsilon}\,\ln\!\frac{2}{\delta} + \frac{8\,d_{\mathrm{vc}}}{\epsilon}\,\ln\!\frac{13}{\epsilon},
\end{equation}
式中，$N$为所需标注样本数，$d_{\mathrm{vc}}$为假设类的VC维度，$\epsilon$为误差容忍度，$\delta$为置信度参数。
\end{theorem}

在EEG应用中，VC维$d_{\mathrm{vc}}$随着通道数、时间长度、频谱维度与模型容量的增加而迅速增大。现代高密度EEG系统可能包含128或256个通道，采样时长可达数分钟，再加上频域特征和深度模型的复杂性，使得$d_{\mathrm{vc}}$达到很高的数值。根据式~\eqref{eq:vc-sample}，这直接推高了获得可靠泛化性能所需的标注样本数$N$。结合多源场景下的功能异构性（定义~\ref{def:hetero}），跨个体可迁移性弱进一步加剧了标注稀缺矛盾，体现出判别式范式在数据效率上的先天不足。

判别式模型还面临无法有效利用无标注数据的根本性问题。在实际应用中，用户在使用BCI系统时会产生大量的EEG信号，但其中只有很小一部分与明确的任务标签相关联。这些无标注数据包含着丰富的神经活动信息和个体特征信息，但判别式方法无法从中学习有用信息，造成了数据资源的巨大浪费。现代深度学习的成功很大程度上依赖于大规模数据的有效利用，而EEG信号领域的数据稀缺问题使得传统深度学习方法难以发挥其应有的优势。

（三）泛化能力的不足

判别式建模的另一个重要局限体现在泛化能力的不足，主要表现在任务特异性、分布敏感性和跨域适应性缺失三个方面。

任务特异性限制源于判别式模型只学习与特定分类任务相关的特征。这些特征往往局限于任务相关的判别性信息，无法捕捉数据的完整结构和内在规律。判别式学习的目标是找到最优的决策边界，这种目标导向的学习方式使得模型倾向于学习任务特定的表面特征，而忽略了数据的深层结构和通用模式。任务特定的特征难以在不同任务间进行有效复用。当面对新的任务或应用场景时，判别式模型往往需要重新采集数据、重新设计特征、重新训练模型，无法利用已有的知识和经验进行快速适应。


分布敏感性是指判别式模型对训练与测试数据分布变化的高度敏感性。在多源EEG环境中，训练与测试阶段常常伴随分布的不一致。判别式范式对分布稳定性高度敏感，其测试风险会随分布偏移而显著放大。为刻画分布偏移对风险的影响，需首先定义全变分距离（Total Variation，TV）。对于概率测度$P$和$Q$，全变分距离定义为
\begin{equation}
\label{eq:tv-def}
\mathrm{TV}(P,Q) \;=\; \tfrac{1}{2}\int |p - q| \,\mathrm{d}\mu \;=\; \sup_{A} |P(A)-Q(A)|,
\end{equation}
式中，$p$和$q$分别为$P$和$Q$关于参考测度$\mu$的密度函数，$A$为可测集合。基于此可建立分布偏移与风险变化的理论关系。

\begin{theorem}[分布偏移下的风险差界\cite{ben2010theory}]
\label{thm:shift}
设训练与测试联合分布分别为$P_{\text{tr}}(X,Y)$与$P_{\text{te}}(X,Y)$，损失$l(\hat{y},y)\in[0,1]$有界。任意模型$f$的训练/测试风险满足
\begin{equation}
\label{eq:tv-bound}
\big| R_{\text{te}}(f) - R_{\text{tr}}(f) \big| \le 2\,\mathrm{TV}\big(P_{\text{tr}}, P_{\text{te}}\big),
\end{equation}
式中，$R_{\text{tr}}(f)=\mathbb{E}_{P_{\text{tr}}}[l(f(X),Y)]$表示训练风险，$R_{\text{te}}(f)=\mathbb{E}_{P_{\text{te}}}[l(f(X),Y)]$表示测试风险。
\end{theorem}

定理~\ref{thm:shift}表明，测试风险可能随着全变分距离$\mathrm{TV}$的增大而显著上升。由于判别式方法缺乏对生成机制的建模，因而难以在分布改变时自适应地校正表示与决策边界。在实际的EEG应用中，分布偏移是一个普遍存在的问题。设备老化会导致硬件特性的缓慢变化；环境条件的改变（如温度、湿度、电磁环境）会影响信号质量；用户状态的变化（如疲劳、情绪、注意力）会改变神经活动模式；长期使用过程中的学习效应也会导致神经模式的演化。这些因素都会造成数据分布的偏移，而判别式模型往往缺乏应对这种偏移的有效机制。

跨域适应性缺失是指判别式模型在面对不同来源或不同域的数据时难以有效迁移的问题。在多源EEG场景中，不同来源的数据往往具有不同的分布特性、结构模式和功能编码方式。判别式模型通常假设所有数据来自相同的分布，无法有效处理这种域间差异。即使采用域适应技术，也往往需要一定数量目标域的标注数据，这又回到了数据稀缺的问题。


\section{生成式建模理论分析}
与直接学习$P(Y\mid X)$不同，生成式建模通过联合刻画$P(X,Z)$与潜在生成机制来学习数据的内在结构与统计规律，从而在表达能力、数据利用与不确定性刻画方面展现优势。本节先给出形式化表达，再围绕其核心优势展开分析与讨论。

\subsection{生成式建模的形式化表达}
生成式建模的核心思想是建立观测数据与潜在变量之间的联合概率模型。设$\mathbf{X}$为观测EEG随机张量，$\mathbf{Z}$为潜在随机变量，$X$和$Z$分别为其对应的观测值，生成式模型学习联合分布：
\begin{equation}
\label{eq:joint}
p_{\theta}(X,Z) \;=\; p_{\theta}(X\mid Z)\,p(Z),
\end{equation}
式中，$p(Z)$为潜在变量的先验分布，$p_{\theta}(X\mid Z)$为给定潜在变量下观测数据的似然函数。真实后验$p_{\theta}(Z\mid X)$往往难以解析计算，故采用变分分布$q_{\phi}(Z\mid X)$进行近似，通过最大化证据下界（Evidence Lower BOund，ELBO）进行优化：
\begin{equation}
\label{eq:elbo}
\log p_{\theta}(X) \;\ge\; \mathbb{E}_{q_{\phi}(Z\mid X)} \big[ \log p_{\theta}(X\mid Z) \big] \;-\; \mathrm{KL}\!\big(q_{\phi}(Z\mid X)\,\|\,p(Z)\big),
\end{equation}
式中，$\theta$为生成模型参数，$\phi$为变分推断参数，$\mathrm{KL}(\cdot\|\cdot)$为Kullback–Leibler散度。右端第一项为重构项，衡量模型重建观测数据的能力；第二项为正则化项，约束潜在表征分布接近先验分布。这种平衡机制使得生成式模型既能够学习数据的复杂模式，又能够保持良好的泛化能力。

与判别式模型只关注输入--输出映射不同，生成式模型通过建模完整的数据分布，能够获得对数据更深层次的理解。这种理解不仅包括数据的表面特征，还包括数据的内在结构、生成机制和统计规律。
EEG信号包含着丰富的时间动态、空间分布和频谱特征，这些特征之间存在复杂的相互关系。生成式模型通过学习这些特征的联合分布，能够捕捉到EEG信号的内在结构和生成机制，为信号分析和应用开发提供更深层次的理解基础。

\subsection{生成式建模的特性分析}

（一）连续表达能力

认知状态往往连续变化，需在表达空间中支持平滑插值与细粒度控制。生成式模型通过连续潜在空间与连续生成映射共同实现了这种连续表达能力。
\begin{theorem}[连续潜在空间的表达优势]
\label{thm:g-exp}
设潜在空间$\mathcal{Z} \subseteq \mathbb{R}^d$为连续空间，生成函数$G: \mathcal{Z} \rightarrow \mathcal{X}$为连续映射，则生成式模型的理论表达能力为不可数无穷。具体地，若$G$为连续映射且$\mathcal{Z}$为连通集，则$|G(\mathcal{Z})| \geq |\mathcal{Z}| = \mathfrak{c}$（连续统基数），从而提供了丰富的表达可能性。
\end{theorem}
\begin{proof}
由于$\mathcal{Z} \subseteq \mathbb{R}^d$且$\mathcal{Z}$为连通集，根据连续映射的性质，连续函数$G$将连通集映射为连通集。当$G$为单射时，$|G(\mathcal{Z})| = |\mathcal{Z}| = \mathfrak{c}$；当$G$不为单射时，由于连续映射的像仍然是连续的，$|G(\mathcal{Z})| \geq \mathfrak{c}_0$（可数无穷）。无论哪种情况，生成式模型的表达能力都远超判别式模型的有限离散输出。更进一步，连续空间的稠密性保证了任意两个不同的点之间都可以通过连续路径连接，这为平滑插值和渐进变化提供了理论基础。
\end{proof}

定理~\ref{thm:g-exp}说明生成式模型天然支持连续表达与潜在空间插值，从而有能力描述意图的渐变与细微差异，这是定理~\ref{thm:disc-cap}所揭示的判别式瓶颈难以实现的。在BCI系统应用中，这意味着系统可以根据用户的神经活动生成全新内容。例如，在艺术创作应用中，系统可以根据用户想象的图像生成对应的视觉内容，不局限于训练数据中的固定类别。在音乐生成应用中，系统可以根据用户的情感状态和音乐偏好生成连续变化的音乐片段，实现个性化的音乐创作。连续表达能力的另一个重要体现是插值和外推能力。生成式模型可以在潜在空间中进行线性插值，生成训练数据中未出现过的中间状态。通过潜在空间的插值，系统可以捕捉和表达这种连续性变化，提供更自然、更符合人类认知特征的交互体验。


（二）自监督与数据效率

生成式建模的第二个优势在于其强大的自监督学习能力。通过自监督方法，生成式模型可以从大量无标注数据中学习有用的表征，而无需获取成本极高的人工标注。
自监督学习的核心思想是设计代理任务，通过这些任务让模型学习数据的内在结构。在EEG信号处理中，常用的代理任务包括信号重构、掩码预测、时序预测和对比学习。信号重构任务要求模型通过编码--解码器架构重构原始信号，从而学习紧致的信号表征。掩码预测任务随机掩盖信号的某些部分，训练模型预测被掩盖的内容，类似于自然语言处理中的BERT模型。时序预测任务根据历史信号预测未来的神经活动模式，利用神经活动的时序相关性来学习有用特征。对比学习通过学习区分相似和不相似的信号片段来获得判别性表征。若自监督预训练可以约束表征类的复杂度，则在相同标注规模下可获得更紧的泛化上界。

\begin{theorem}[预训练降低复杂度的泛化改进\cite{mohri2018foundations}]
\label{thm:ss-generalize}
令$\mathcal{F}\circ\mathcal{H}$为下游分类器族与表征族的复合类。设利普希茨有界损失的经验Rademacher复杂度为$\mathfrak{R}_N(\mathcal{F}\circ\mathcal{H})$。若自监督预训练得到的表征族$\mathcal{H}_{\text{pre}}$满足
\begin{equation}
\label{eq:rademacher-reduce}
\mathfrak{R}_N(\mathcal{F}\circ\mathcal{H}_{\text{pre}}) \;\le\; \gamma \,\mathfrak{R}_N(\mathcal{F}\circ\mathcal{H}), \quad \gamma\in(0,1),
\end{equation}
则在同一$N_l$下，基于Rademacher复杂度的标准泛化上界按系数$\gamma$收紧；若以$\mathfrak{R}_N\!\propto\! 1/\sqrt{N}$近似，则达成相同上界所需标注量近似按$N_l^{\text{pre}} \approx \gamma^{2} N_l$缩减。
\end{theorem}

定理~\ref{thm:ss-generalize}表明，当自监督预训练成功降低表征类的Rademacher复杂度时（即$\gamma<1$），可在固定标注样本数下实现更紧的泛化上界，或等价地，在目标泛化性能下显著减少所需的标注样本数。

（三）不确定性建模能力

生成式建模的第三个优势体现在其对不确定性的显式建模能力。EEG信号具有内在的随机性和变异性，相同的认知状态可能对应多种不同的神经活动模式。这种不确定性来源于多个方面。首先是观测不确定性，由于设备噪声、环境干扰、生理伪影等因素的存在，观测到的信号与真实的神经活动之间存在差异。其次是认知不确定性，相同的认知状态可能通过不同的神经活动模式来实现，这反映了大脑功能的冗余性和可塑性。最后是个体不确定性，不同个体在执行相同认知任务时可能表现出不同的神经活动模式，这反映了个体间的生理和认知差异。

% EEG具有观测、认知与个体三重不确定性。生成式模型以概率分布直接刻画这些不确定性，为风险评估与可靠决策提供基础。
\begin{definition}[EEG信号中的不确定性]
\label{def:uncertainty}
设$X$为观测EEG信号，$Z$为潜在神经活动因子，$c \in \mathcal{C}$为认知状态标签，$s \in \mathcal{S}$为个体身份标识，$f(\cdot)$为确定性生成函数，$\epsilon$为观测噪声，$\Sigma$为噪声协方差矩阵。EEG信号中的不确定性主要包括以下三个方面：
\begin{equation}
\label{eq:uncertainty}
\left\{
\begin{aligned}
\text{观测不确定性：}&\quad X \;=\; f(Z) + \epsilon,\;\; \epsilon \sim \mathcal{N}(0,\Sigma);\\
\text{认知不确定性：}&\quad p(Z\mid c)\;\text{为多峰分布，}\; \mathbb{H}[p(Z\mid c)] > \tau;\\
\text{个体不确定性：}&\quad p(X\mid c,s) \;\neq\; p(X\mid c,s'),\;\; \forall\, s\neq s' \in \mathcal{S},
\end{aligned}
\right.
\end{equation}
式中，$\mathbb{H}[\cdot]$表示微分熵，$\tau > 0$为熵阈值，$p(Z\mid c)$表示给定认知状态$c$下潜在因子$Z$的条件分布，$p(X\mid c,s)$表示给定认知状态$c$和个体$s$下观测信号$X$的条件分布。
\end{definition}

如定义~\ref{def:uncertainty}所示，EEG信号中的三重不确定性导致了复杂的一对多映射关系，即相同的认知状态可对应多种不同的神经活动模式。判别式建模通常假设确定性的输入--输出映射，难以有效处理这种内在的随机性。相比之下，生成式模型通过概率分布建模这些不确定性，能够更准确地刻画神经活动的随机本质，为风险评估、置信度估计和鲁棒决策提供理论基础。

（四）通用表征与迁移

% 跨个体/跨任务共享的神经机制意味着潜在语义空间可对齐。若源/目标潜在先验与条件生成机制相容，则目标域过量风险可控。
生成式建模还具有强大的语义理解能力。通过学习从低维潜在变量到高维观测数据的映射，生成式模型能够发现数据中的语义结构。潜在变量往往对应于数据的语义属性，如认知状态、情感状态、注意水平等。通过操控潜在变量，可以生成具有特定语义特征的数据，这为理解和控制数据的语义内容提供了途径。
知识迁移能力源于生成式模型学习到的通用表征。这些表征不是针对特定任务的，而是反映数据的内在结构和规律。迁移学习的成功依赖于源任务和目标任务之间的相似性，生成式模型通过学习数据的生成过程，能够发现不同任务间的共同结构。

\begin{theorem}[潜在对齐下的跨任务迁移界]
\label{thm:transfer}
设源/目标任务潜在先验为$p_s(z)$、$p_t(z)$，存在可测变换$T:\mathcal{Z}_s\to\mathcal{Z}_t$使推送测度$T_{\#}p_s$满足
\begin{equation}
\label{eq:latent-align}
\mathrm{KL}\!\big(p_t(z)\,\|\,T_{\#}p_s(z)\big)\;\le\;\varepsilon,
\end{equation}
且条件似然的差异满足$\sup_{z}\mathrm{TV}\!\big(p_t(x\mid z),\,p_s(x\mid T^{-1}(z))\big)\le \eta$。则对任一有界损失，下游解码器在目标域的过量风险相对源域受控于$O(\varepsilon+\eta)$。
\end{theorem}
\begin{proof}
记源域和目标域的联合分布分别为$P_s(x,z) = p_s(x|z)p_s(z)$和$P_t(x,z) = p_t(x|z)p_t(z)$。通过变换$T$对齐后，源域分布变为$P_s^T(x,z) = p_s(x|T^{-1}(z))T_{\#}p_s(z)$。

根据全变分距离的性质和积分形式有：
\begin{equation}
\begin{aligned}
\mathrm{TV}(P_t, P_s^T) &= \frac{1}{2}\int |p_t(x,z) - p_s(x|T^{-1}(z))p_{T_{\#}s}(z)| dxdz\\
&\le \frac{1}{2}\int p_t(z)|p_t(x|z) - p_s(x|T^{-1}(z))| dxdz\\
&\quad + \frac{1}{2}\int |p_t(z) - p_{T_{\#}s}(z)| dz\\
&\le \eta + \varepsilon
\end{aligned}
\end{equation}
由分布偏移下的风险差界（定理~\ref{thm:shift}），有界损失下的过量风险满足：
\begin{equation}
|R_t(f) - R_s^T(f)| \le 2\,\mathrm{TV}(P_t, P_s^T) \le 2(\varepsilon + \eta)
\end{equation}
因此过量风险受控于$O(\varepsilon+\eta)$。
\end{proof}

定理~\ref{thm:transfer}刻画了生成式模型跨任务迁移的理论保证。其中变换$T$表示源域与目标域潜在空间之间的对齐映射，KL散度项$\varepsilon$量化了经过对齐后的源域先验分布与目标域先验分布之间的差异，全变分距离项$\eta$度量了条件生成机制在对齐后的潜在空间中的一致性程度。当潜在空间能够成功对齐（即$\varepsilon$较小）且生成机制具有良好的可迁移性（即$\eta$较小）时，目标域的过量风险相对于源域的增长可以被有效控制在$O(\varepsilon+\eta)$的量级。这一理论结果为生成式模型在跨被试、跨设备和跨任务场景下的迁移学习提供了理论基础，同时也回应了定义~\ref{def:hetero}中关于多源EEG环境下语义异构和结构异构问题的讨论。


\section{双层级生成式建模框架构建}
基于生成式建模的理论优势分析，本文建立了信号级--表征级的双层级建模框架，契合神经信息处理的层次性特征。该框架的设计灵感来源于大脑信息处理的层次化组织结构。从底层的神经元活动到高层的认知功能，每个层级都承担着特定的信息处理任务。双层级生成式框架通过分层建模的方式，在底层聚焦信号质量与统计一致性的改善，在上层学习具有强泛化能力的通用表征。这种分层递进的设计充分发挥了生成式建模在连续表达、自监督学习和不确定性建模方面的理论优势，为复杂的多源EEG应用场景提供了系统性的解决方案。


\subsection{信号级生成建模理论基础}

在EEG应用中，信号质量问题是一个普遍存在的挑战。原始采集的EEG信号往往包含各种类型的噪声和伪影，如眼电伪影、肌电干扰、心电伪影、工频噪声、基线漂移等。这些干扰成分的幅度往往远大于真实的神经信号，严重影响后续分析的准确性。传统的信号处理方法如滤波、独立成分分析等虽然能够在一定程度上改善信号质量，但往往需要人工调参，且对不同类型的干扰效果差异较大。生成式建模为信号质量改善提供了新的思路。
信号级生成式建模的目标在于建模观测分布并提升信号可用性，包括去噪与去伪影、缺失值插值、通道超分重建、统计标准化等多个方面，从而为上层表征学习提供稳定分布与高信噪比的输入。通过学习清洁信号的分布特征，生成式模型能够自动区分真实的神经信号和各种干扰成分，实现智能化的信号增强。

针对EEG信号的去噪与重建需求，$\beta$-VAE提供了一个有效的理论框架\cite{higgins2017beta}。该方法在标准VAE的ELBO基础上引入权重参数$\beta$，通过调节重构保真度与表征压缩之间的率--失真（Rate-Distortion）权衡来实现信号增强。标准VAE的ELBO本质上对应于最小化重构失真与编码率的加权和，其中编码率可近似为$R\!\approx\!\mathbb{E}_{q_{\phi}(z\mid x)}[\mathrm{KL}(q_{\phi}(z\mid x)\,\|\,p(z))]$。当$\beta>1$时，模型更强调表征压缩与解耦合，鼓励潜在变量$z$承载更简洁、近似独立的语义因子；当$\beta<1$时，模型更追求重构保真度。其优化目标为：
\begin{equation}
\label{eq:loss-signal}
\mathcal{L}_{\text{sig}}
\;=\;
\mathbb{E}_{q_{\phi}(z\mid x)} \big[ \log p_{\theta}(x\mid z) \big]
\;-\;
\beta\, \mathrm{KL}\big(q_{\phi}(z\mid x)\,\|\,p(z)\big),
\end{equation}
式中，$x$表示输入信号段，$z$表示潜在变量，$q_{\phi}(z\mid x)$为编码器参数化的近似后验分布，$p_{\theta}(x\mid z)$为解码器参数化的条件似然，$p(z)$为潜在变量的先验分布，$\phi$和$\theta$分别表示编码器和解码器的参数，$\beta \geq 0$为权衡参数。在率-失真理论的视角下，参数$\beta$控制着编码码率，进而影响模型对噪声的吸收能力与有效信息的保留程度。较大的$\beta$值会促使模型学习更压缩的表征，自然地过滤掉噪声和不重要的细节，但可能会损失一些有用信息；较小的$\beta$值则允许模型保留更多细节，但可能会保留噪声成分。

针对EEG信号的特殊性质，可以在基本框架基础上引入稳健似然函数与物理先验知识。对于时域重构任务，可以使用带外噪声抑制的Laplace/Gaussian混合似然，这种混合分布能够更好地处理信号中的突发性噪声和异常值。对于频域重建任务，可以使用对数功率谱的高斯似然，这种设计符合EEG信号功率谱的统计特性。当存在跨设备域偏移时，可以在潜在空间施加域对齐约束，如最大均值差异约束或对抗域对齐。也可以通过显式观测算子正则化$\|X - \mathcal{T}_s(A_s \tilde{S})\|$将体积传导一致性纳入优化目标，其中$\mathcal{T}_s$表示变换算子，$A_s$表示观测算子，$\tilde{S}$表示重建的源信号。这些设计使信号级模型既能作为强大的去噪与插值器，又能作为上层表征学习的标准化前端，为整个生成式建模框架提供高质量的数据基础。


\begin{figure}[htbp] % use float package if you want it here
  \centering
  \includegraphics[width=0.8\textwidth]{C2-1sig.pdf}
  \caption{信号级生成建模理论基础和具体实现}
  \label{fig:C2-1}
\end{figure}

如图~\ref{fig:C2-1}所示，本文在信号级生成建模技术体系中重点涉及空间增强MASER和通道选择CogSSM两种关键方法。其中，MASER专注于从低密度到高密度EEG信号的空间超分辨率重建，通过状态空间建模技术实现信号的空间增强；CogSSM则从通道选择的角度实现信号约简，在保持关键神经信息完整性的前提下，通过协同引导学习策略智能化地识别最优电极配置。这两种方法分别沿着增强和约简两个维度构建了协同互补的技术方案，使信号级模型既能作为强大的去噪与插值器，又能作为上层表征学习的标准化前端，为整个生成式建模框架提供高质量的数据基础。



\subsection{表征级生成建模理论基础}

表征级的目标在于学习映射函数$h:\mathcal{X}\to\mathcal{H}$，使得语义相近的样本在表征空间$\mathcal{H}$中形成紧密聚类，同时实现跨源对齐、时序一致性和信息可重构性等多重目标。表征级建模是连接底层信号和高层应用的关键桥梁，其学习到的表征应该既保留重要的神经信息，又具有良好的泛化和迁移能力。

大脑的信息处理具有层次化和分布式的特点，这种层次性在EEG信号的频谱、时间和空间维度上都有体现。低层表征主要编码信号的基础物理特征，如单通道幅值、局部频谱能量等；中层表征捕捉时空模式，如事件相关电位的峰值潜伏期、跨通道相位同步等；高层表征则对应复杂的认知功能模式，如注意状态、工作记忆负荷等抽象概念。表征级建模的核心挑战在于如何在保持层次化结构的同时，确保不同抽象级别的表征都具有可解释性和可迁移性。
针对EEG信号的多源异构特性，表征级建模需要统筹考虑重构保真度、跨源一致性、时序连续性和语义判别性等多个相互制约的目标。一种有效的策略是采用复合自监督学习框架，将多个互补的学习目标耦合在统一的优化过程中，例如：
\begin{equation}
\label{eq:loss-repr}
\begin{aligned}
\mathcal{L}_{\text{repr}}
=\;
\mathbb{E}_{x \sim \mathcal{D}} \Big[
&\underbrace{-\log p_{\theta}\!\big(x \mid h(x)\big)}_{\text{重构}}\;
+\;
\lambda_c\, \underbrace{\Big(\!-\log \frac{\exp(\mathrm{sim}(h(x),h(x^+))/\tau)}{\sum_{x^-}\exp(\mathrm{sim}(h(x),h(x^-))/\tau)}\Big)}_{\text{对比}}\\
&+\;
\lambda_p\, \underbrace{\big\|\, h(x_{t}) - f_{\text{pred}}\!\big(h(x_{t-k:t-1})\big)\,\big\|_2^2}_{\text{时序预测}}
\;+\;
\lambda_a\, \underbrace{\mathrm{MMD}\!\big(h(x_{s_1}),\, h(x_{s_2})\big)}_{\text{跨源对齐}}
\Big],
\end{aligned}
\end{equation}
式中，$\mathrm{sim}(\cdot,\cdot)$采用余弦相似度度量，$\tau$为温度参数控制对比学习的锐度，$x^+$表示同源增强样本或语义相关的时间邻域，$x^-$来自不同被试、设备或任务的负样本。最大均值差异（Maximum Mean Discrepancy，MMD）衡量两组数据在某个核诱导的特征空间中均值嵌入的距离，从而比较它们是否来自同一分布。实际计算中用核函数把样本对相互比较，得到无偏或有偏的经验估计来近似这个距离。

式~\eqref{eq:loss-repr}中各项的协同作用机制如下：重构项$-\log p_{\theta}(x \mid h(x))$确保表征空间保留足够的信息量，避免过度压缩导致的信息丢失，同时提供生成式可逆性以支持数据合成和插值；对比项通过拉近语义相似样本、推远无关样本的方式增强表征的判别性，有效压缩任务无关的个体差异和噪声变异；时序预测项$f_{\text{pred}}$基于历史窗口$x_{t-k:t-1}$预测当前表征$h(x_t)$，编码神经动力学的时序依赖性和因果结构；跨源对齐项通过最小化不同数据源表征分布的MMD距离，缓解设备差异、个体变异等因素造成的域偏移问题。为进一步促进表征的解耦性和可解释性，可在$h(x)$上施加结构化约束，如正交性正则化、独立性先验，或针对不同频带的子表征采用块对角结构惩罚。

\begin{figure}[htbp] % use float package if you want it here
  \centering
  \includegraphics[width=0.6\textwidth]{C2-2repr.pdf}
  \caption{表征级生成建模理论基础和具体实现}
  \label{fig:C2-2}
\end{figure}


如图~\ref{fig:C2-2}所示，本文在表征级生成建模中采用基于去噪掩码自编码器的DMAE-EEG方法。该方法通过将神经科学先验知识与深度学习技术有机整合，构建了统一的自监督学习框架。核心创新包括脑区拓扑异质划分方法，将不规则的电极分布投影为稳定的脑区块--补丁图式化表征；去噪伪标签生成器通过重建型去噪任务生成稳健的教师信号；非对称掩码自编码器架构利用Transformer的自注意力机制捕获长程时空依赖关系，实现了从大规模无标注数据中学习具有强泛化能力的通用表征。
表征级建模的输出具有多重功能属性。在下游应用中，学习到的表征$h(x)$可直接作为分类器、回归器等判别模型的输入特征，提供比原始信号更紧致、更具判别性的信息编码；在生成任务中，表征可作为条件变量指导高层的语义生成过程，如基于认知状态的信号合成、跨被试的模式迁移等。



通过信号级与表征级的协同设计，整个生成式建模框架实现了从底层信号质量提升到高层语义理解的端到端优化。信号级模块专注于噪声抑制、伪影去除和数据标准化，为上层分析提供高信噪比的输入；表征级模块则在清洁信号基础上学习层次化的语义表征，兼顾判别性和生成性需求。这种分层递进的架构设计不仅充分发挥了生成式建模在无监督学习和跨域泛化方面的优势，还为复杂的多源EEG应用场景提供了实用有效的建模范式。


% \section{本章小结}
% 本章从理论角度系统构建了多源EEG信号生成式建模的数学基础，为后续的方法设计和实验验证奠定了坚实的理论根基。通过张量化表示、概率空间描述和信息结构定义，本章首先给出了多源EEG信号的完整形式化表示，明确了其在分布、结构和功能三个层面的异构特性。这种形式化表示不仅为理解多源数据的复杂性提供了统一框架，也为后续建模方法的设计指明了方向。
% 在此基础上，本章深入分析了传统判别式建模范式在处理多源EEG数据时面临的根本性局限。通过严格的理论分析，揭示了判别式方法在表达能力、数据效率和泛化能力三个关键方面的不足。表达能力的有界性限制了系统的功能多样性，数据效率的低下制约了实际部署的可行性，泛化能力的不足影响了跨域应用的效果。这些理论分析不仅解释了现有EEG-BCI系统的局限性，也为生成式建模方法的引入提供了充分的理论依据。
% 相比之下，生成式建模在连续表达能力、自监督学习、不确定性建模和通用表征学习等方面展现出的显著优势，为克服判别式方法的局限提供了有效途径。连续潜在空间的无限表达能力突破了离散标签的束缚，自监督学习策略大幅降低了对标注数据的依赖，显式的不确定性建模提供了更可靠的决策基础，通用表征的学习能力增强了跨任务和跨域的知识迁移。这些优势的理论分析为生成式方法在EEG领域的应用提供了坚实的理论支撑。
% 基于这些理论分析，本章最终提出了信号级-表征级的双层级生成式建模框架。信号级专注于提升信号质量和统计一致性，通过去噪、插值、标准化等操作为后续处理提供高质量的数据基础；表征级致力于学习通用的、可迁移的特征表示，通过复合自监督目标实现重构、对比、时序建模和跨源对齐的统一优化。这种分层设计充分体现了神经信息处理的层次性特征，每一层都有明确的功能定位和优化目标，层间通过信息流动实现有机衔接。
% 所提出的理论框架不仅为后续章节的具体方法设计提供了指导原则，也为整个研究领域的发展指明了方向。通过将生成式建模的理论优势与EEG信号的特殊性质相结合，这一框架有望推动EEG-BCI技术从实验室研究向实际应用的转化，为构建更智能、更鲁棒、更实用的脑机接口系统奠定重要的理论基础。

